\section{Results}

\subsection{Historical Failure Atlas}

The registry contains 2,209 V1--V22 rows, including 1,637 paired Delta ARI rows
grouped into 431 dataset/protocol/readout units. The paired atlas contains 194
material-positive, 680 material-negative, and 763 observed-small rows at
$\delta=0.03$. The signs and magnitudes vary across intervention families and
protocols, but the registry does not provide one common structural-opportunity
endpoint. We therefore report these values as an observational atlas rather than
as a pooled causal estimate.

\begin{table}[t]
\centering
\caption{Counts in the frozen retrospective atlas. Rows are repeated records,
not independent data sets.}
\label{tab:atlas}
\begin{tabular}{lr}
\toprule
Quantity & Value \\
\midrule
V1--V22 registry rows & 2209 \\
Paired $\Delta$ARI rows & 1637 \\
Dataset/protocol/readout units & 431 \\
Material positive rows & 194 \\
Material negative rows & 680 \\
Observed-small rows & 763 \\
\bottomrule
\end{tabular}

\end{table}

\begin{figure}[t]
\centering
\includegraphics[width=0.98\linewidth]{V25_Figure1_failure_atlas}
\caption{V1--V22 Failure Atlas. The display is observational and stratified by
version/family; it is not a pooled causal estimate.}
\label{fig:atlas}
\end{figure}

\subsection{Matched selection effects}

The audited E1 panels pass the full matching and label-isolation contract. The
per-seed primary endpoint is $S_{\mathrm{full\_ARI}}$; the confirmation dataset
means $S_d=\operatorname{mean}_s(S_{\mathrm{full\_ARI}})$ are shown in
Table~\ref{tab:e1} and Figure~\ref{fig:e1}.
The primary selection effect changes sign across the three confirmation data
sets. The generic intervention effect is also heterogeneous: it is positive for
Campbell, observed-small for \texttt{hate\_speech}, and inconclusive for Baron
Human. This decomposition separates generic intervention from the topology
policy increment; it does not prove either component universal.

\begin{table}[t]
\centering
\caption{Confirmation means for the matched V21 decomposition. Each row is a
dataset mean over three repeated seeds.}
\label{tab:e1}
\begin{tabular}{lrrl}
\toprule
Dataset & $I_d$ & $S_d$ & $S_d$ state \\
\midrule
Baron Human & -0.011463 & +0.044617 & Positive \\
Campbell & +0.132841 & -0.065332 & Negative \\
hate\_speech & -0.002671 & -0.033410 & Negative \\
\bottomrule
\end{tabular}

\end{table}

\begin{figure}[t]
\centering
\includegraphics[width=0.78\linewidth]{V25_Figure3_e1_selectivity}
\caption{Generic intervention effect $I_d$ versus topology selectivity effect
$S_d$. The signs are conditional case-study evidence, not a population claim.}
\label{fig:e1}
\end{figure}

The pilot gate passed because at least two of three pilot datasets showed a
seed-stable material effect in either $I_d$ or $S_d$; signs were not required to
agree. Pilot and confirmation are not pooled: the confirmation panel alone
supplies the three case-study values above. The result is therefore evidence for
conditional utility and heterogeneity, not a reason to search for a universally
beneficial selector.

\subsection{Localization diagnostics}

E2 produces dataset-by-seed semantic summaries and gradient/one-step records
under the frozen N/R/T protocol. These artifacts permit inspection of selected
feature semantics and optimizer-direction changes, but they do not produce a
stable cross-dataset objective-conflict rule. Figure~\ref{fig:diagnostics}
shows the diagnostic geometry without fitting a cross-dataset causal model.

\begin{figure}[t]
\centering
\includegraphics[width=0.98\linewidth]{V25_Figure4_diagnostics}
\caption{E2 diagnostic geometry. Gradient cosines and one-step/full effects are
descriptive records; they are not a universal objective-compatibility law.}
\label{fig:diagnostics}
\end{figure}

The E3 replay gate found zero artifact-complete candidate rows, so no replay was
run. The separate V23 boundary rows are retained only to show why local and
global readouts must be measured separately; they do not support a universal
conversion law.

\begin{figure}[t]
\centering
\includegraphics[width=0.78\linewidth]{V25_Figure5_local_global_boundary}
\caption{V23 local/global boundary evidence. These rows are kept outside the
V1--V22 atlas and are not a universal theorem.}
\label{fig:boundary}
\end{figure}

\subsection{Independent validation status}

The holdout has no evaluable endpoint: zero of six panels completed under the
frozen protocol, and the recorded cause is the dense decoder/Adam resource
boundary. Independent validation is therefore not established. We do not use
this incomplete compute to falsify the selection claim or report a negative
performance result.
